\paragraph{Bounds obtained from all possible coarse splittings.}
The available bundle data give the further bounds
\[
 0\le h^0(K_m)\le 2m+2,\qquad
 4m+1\le h^{1,1}(X_m)\le6m+3,\qquad
 4m\le h^{1,2}(X_m)\le6m+2.
\]
Here is a complete derivation. Put $S=2m-1$, and write
\[
 V_c=\rho_*E_m=\OO(a)\oplus\OO(S-a),\qquad
 G_c=\rho_*G_m=\OO(b)\oplus\OO(S-b).
\]
The already proved nonnegative splitting implies $0\le a,b\le S$.
Reordering each pair gives $a,b\le S/2$.
Taking invariants defines a natural injective morphism
\[
 \rho_*\mathcal Hom(E_m,G_m)\longrightarrow
 \mathcal Hom(V_c,G_c).
\]
It is an isomorphism away from the two stack points, and its source is
torsion free, which proves injectivity. In the invariant local bases,
the map from the nontrivial-character source line to the trivial-character
target line has to vanish at the stack point; the other entries have no
additional condition. Thus there is one elementary condition at each
finite point, consistently with the source degree $-2$ in
\eqref{invII:eq:E.7}.

The four contributions to $h^0(\mathcal Hom(V_c,G_c))$ are
\[
 \max\{b-a+1,0\},\quad S-a-b+1,\quad
 \max\{a+b-S+1,0\},\quad\max\{a-b+1,0\}.
\]
If $a+b<S$, the third is zero and their sum is at most
$S-a-b+1+|a-b|+2\le S+3$. If $a+b=S$, the two inequalities
$a,b\le S/2$ imply $a=b=S/2$, and the sum is four, also at most
$S+3$. Since $K_m$ is a subbundle of
$\rho_*\mathcal Hom(E_m,G_m)$, this proves $h^0(K_m)\le S+3=2m+2$.
Substitution in \eqref{invII:eq:E.5} proves the two Hodge bounds.
In particular $0\le h^0(K_{12})\le26$. No balanced splitting has been
assumed, and these bounds are not asserted to be sharp.

\paragraph{Why the numerical constraints do not determine the answer.}
Set $D=\OO^{\oplus2}\oplus\OO(-1)^{\oplus2}$ and
$L=\OO(2m)$ on $\PP^1$, with homogeneous coordinates $[X:Y]$.
The two morphisms
\[
 \phi_A=(X^{2m},Y^{2m},0,0),\qquad
 \phi_B=(0,0,X^{2m+1},Y^{2m+1}):D\longrightarrow L
\]
are surjective, because their two nonzero entries have no common
projective zero. The degree-$2m$ pencil sequence, together with the
unused summands, gives
\[
 \ker\phi_A=\OO(-2m)\oplus\OO(-1)^{\oplus2}.
\]
The degree-$(2m+1)$ pencil sequence tensored by $\OO(-1)$ gives
\[
 \ker\phi_B=\OO^{\oplus2}\oplus\OO(-2m-2).
\]
Both maps therefore have rank-four source of degree $-2$, target of
degree $2m$, and cokernel of length zero. Their rank-three kernels
both have degree $-2m-2$, but their section dimensions are zero and
two. For $m=12$ the splitting triples are $(-24,-1,-1)$ and
$(0,0,-26)$. These exact vector-bundle examples show that the known
rank, degree and cokernel length alone cannot determine the kernel.
We do not assert that either morphism is the cup map of an admissible
comb period section, or that different such geometric sections have
different Hodge numbers.

The retained invariant homology and cohomology classes do not remove
this missing information. They supply sections of $E_m$ and $G_m$,
and their tensor belongs to $E_m\otimes G_m$. The cotangent kernel
instead involves $E_m^*\otimes G_m$; there is no specified identification
of $E_m$ with $E_m^*$. Likewise the exact tangent calculation uses
a different cup map. An actual admissible period section, its transition
matrices and its derivative are required to decide the remaining
kernel. The existence and smoothing construction supplies an open set
of such sections, rather than coefficient data selecting a unique
one. The estimates and the algebraic examples above do not substitute
for solving that geometric choice and its cup-map calculation.
